% Book-specific preamble, \input by en/main.tex and ja/main.tex (the
% repository root is on TEXINPUTS, so this also works in chapter-only
% builds).  Put this book's own definitions here — listing languages
% (\lstdefinelanguage), shorthand macros, extra packages — instead of
% editing format/lnotes.cls, which is shared by every book.

% PTX fragments (Lessons 4, 10, 11): a minimal listings language.
% Opcodes and state-space / type suffixes are keywords; %-registers are
% plain text; "//" starts a comment.
\lstdefinelanguage{PTX}{
  morekeywords={ld,st,add,sub,mul,mad,mov,setp,bra,ret,cvt,shl,shr,
                global,shared,local,param,u32,s32,f32,b32,pred,
                reg,entry,visible,version,target},
  sensitive=true,
  morecomment=[l]{//},
}

% Floating-point operations as a siunitx unit (Lessons 4, 8, 12):
% \SI{1}{\tera\flop\per\second} prints 1 TFLOP/s.
\DeclareSIUnit{\flop}{FLOP}

% Units: \per prints as a slash (GB/s, TFLOP/s), not as a negative power.
% Japanese edition: a range prints as "20--50 µs" instead of the English
% word "to" (the same setting as the cs6290 and cs6200 books, there per call).
\sisetup{per-mode=symbol}
\iflangja{\sisetup{range-phrase={--}, range-units=single}}{}
